% 本文件由 assemble.py 按 _work/plans/ch02.json 逐字组装，请勿手改（改 plan 或 blocks 后重新生成）。
\chapter{相互作用——力}
\label{ch:02}

\section{重力弹力摩擦力}
% ---- block 071-01 (原稿第71页) kind=summary ----
% 页面左侧边缘露出的「水/飞/火/汽车/水平力的合成与分解/圆/水平/单摆」等字属相邻页边缘，未转录
\textbf{（页边）相互作用}

\textbf{重力\quad 弹力\quad 摩擦力}
\begin{itemize}
\item 弹力的分析与计算
  \begin{itemize}
  \item 用胡克定律
  \item 力的平衡条件
  \item 牛顿第二定律
  \end{itemize}
\item 摩擦力的有无与方向判断：假设法\quad 状态法\quad 转换法
\item 摩擦力大小的计算
  \begin{itemize}
  \item 滑摩：$f_{\text{滑}}=\mu N$\quad $\pm$\quad $\mu mg$ % CHECK: 「=μN」与「μmg」之间的符号像「±」（也可能是「=」），按原样转录
    \\ 或由平衡条件、牛二求
  \item 静摩：一般 $f_{\text{静max}}=\mu N$\quad 有时 $f_{\text{静max}}>\mu N$
    \\ 或由平衡条件、牛二求.
  \end{itemize}
\item 摩擦力突变问题：快带慢\quad 慢阻快\quad $+$\quad 快后慢\unclear{断}
  \\ $\leftarrow$ $v$的关系变了\quad （手写箭头由「$v$的关系变了」指向「摩擦力突变问题」）
\end{itemize}


\section{力的合成与分解}
% ---- block 071-02 (原稿第71页) kind=summary ----
\textbf{力的合成与分解.}\quad \textbf{共点力}

\begin{tabular}{@{}lll@{}}
力的合成 & 等效替代法 & 平行四边形法则 \\
 & 平移 & 三角形法则 \\
力的分解 & 正交分解. & \\
\end{tabular}


\section{活结死结与动杆定杆}
% ---- block 071-03 (原稿第71页) kind=knowledge ----
\textbf{活结与死结、动杆与定杆模型}

活结两侧绳张力相等、死结绳两侧力不一定相等、动杆弹力必沿杆，定杆弹力方向任意


\section{受力分析与共点力平衡}
% ---- block 071-04 (原稿第71页) kind=summary ----
\textbf{受力分析、共点力的平衡}
\begin{itemize}
\item \unclear{顺序}：重弹摩场、合成法\quad 正交分解法.
\item 静态平衡
\item 临界：极限分析法、数学分析法（二次函数、三角函数、不等式）\quad 图解法.
\item 动态平衡.
  \begin{itemize}
  \item 解析法
  \item 图解法
  \item 相似三角形法.
  \item 正交分解法
  \item 力辅助圆 % 页面底部被照片截断，此行只拍到一半
  \end{itemize}
\end{itemize}

% ---- block 072-01 (原稿第72页) kind=summary ----
% 本页为横置页（已转正）；左缘倾斜的小字、红色封面等属相邻页/杂物，未转录
\textbf{力学\quad 共点力平衡}

必然\quad 所有力延长线交于一点

\medskip
\textbf{静态平衡}\quad （求某个力的值）
% 手写箭头：「在一条线上」右侧有箭头指向「等大反向」；另有一斜细箭头自此处指向下方「三力」一行；「力三角形闭合」「力的多边形闭合」与「建系…」之间有一个空心粗箭头
\begin{itemize}
\item 二力\quad 在一条线上 $\to$ 等大反向
\item 三力\quad 力三角形闭合 $\Rightarrow$ 建系\ $\sum F_x=0$\ \ $\sum F_y=0$
\item 多力\quad 力的多边形闭合 $\Rightarrow$ 建系\ $\sum F_x=0$\ \ $\sum F_y=0$
\item （页面右侧另有一列，自上而下）有直角用勾股定理 / 非直角三角形 / 正弦定理 / 余弦定理
\end{itemize}

\medskip
\textbf{动态平衡}\quad （求力的变化）
\begin{itemize}
\item 二力\quad 等大反向
\item 三力\quad 力三角形\quad 受力分析：
  \begin{itemize}
  \item 旋转\quad 2个力方向不变\quad 、延长，
  \item 相似\quad 有绳子/杆\quad 非特殊三角形\quad 相似\ 正\unclear{余}弦定理
  \item 力辅助圆\quad 有一角不变\quad 有\unclear{个}力不变
  \item 建系分解.
  \end{itemize}
\item 多力 % 位置在「动态平衡」之下、与二力三力同列
  \begin{itemize}
  \item 4个力\quad 有摩擦\quad 有支持力\quad 用摩擦角
  \item 没有摩擦角\quad 正交分解到$X$、$Y$方向\,方程
  \end{itemize}
\end{itemize}

% ---- block 072-02 (原稿第72页) kind=method ----
\textbf{力学两确定.}
\begin{itemize}
\item 平衡还是非平衡
\item 研究对象是谁
\end{itemize}


\section{共点力平衡·临界极值}
% ---- block 086-02 (原稿第86页) kind=problem ----
\begin{problem}[1、]
\notefig[0.3]{figures/p086_d.png}
\unclear{$\mu$}$=0.125$\quad 系统静止，求 $m_B$ 范围
% CHECK: 手写像 "M=0.125"，按物理应为动摩擦因数 μ=0.125（斜面 37°，A、B 由绳经斜面顶端定滑轮相连，A 在 B 上）；答案中的 M 即 m_B，m 为 m_A
\begin{solution}
\[ \begin{cases} m_{B\,\max}\text{ 时 }f_B\text{ 向上}\\[2pt] m_{B\,\min}\text{ 时 }f_B\text{ 向下}\end{cases}\qquad \frac37m\le M\le\frac95m \]
\end{solution}
\end{problem}


\section{共点力平衡}
% ---- block 114-03 (原稿第114页) kind=problem ----
\begin{problem}[3]
% 图p114_d内含标注：水平线上的T、θ、mg、F方向等（受力示意）；图p114_e：F、θ、mg、T 的力三角形
\notefig[0.35]{figures/p114_d.png}

$F=\rho SV^2\sin^2\theta$\quad 求 $V_{\min}$
\begin{solution}
\[mg=F\cos\theta\]
\[\therefore\ V=\sqrt{\frac{mg}{\rho S\sin^2\theta\cos\theta}}\qquad \tan\theta=\frac{\sqrt{2}}{2}\ \text{时}\ V_{\min}\] % CHECK: 若 F∝sin^2θ，使 sin^2θ cosθ 最大应得 tanθ=√2；原稿写作 √2/2（拍照看似 √2 over 2）

\notefig[0.25]{figures/p114_e.png}
\end{solution}
\end{problem}

